\vspace{-0.2cm}
\reviewsection{The READ Turned Support Into Adult Accountability}

\vspace{-0.2cm}
The Module 2 readings give me a disciplined way to move from principle to practice. Dunlap et al. organize individualized positive behavior support through teaming and goal setting, data collection, functional behavioral assessment, intervention planning, and progress monitoring \citep{dunlap2019ptr}. The Prevent--Teach--Reinforce structure asks adults to identify conditions they can change, teach a useful communication or participation route, make that route effective for the learner, and revise the support using evidence. The most important shift is that support data must include what adults and environments do. If directions were inaccessible, a promised option was unavailable, processing time was interrupted, or a communication route was ignored, a record of student behavior alone is incomplete. A strategy respects dignity when the goal is meaningful in the learner's life; consent when the learner can influence, refuse, pause, and return; and neurodiversity when progress is not defined as appearing less disabled.

My teaching placement gave me a practical example of that responsibility. Participation changed when I reduced language demands, modeled technology steps, and allowed flexible ways of engaging with an activity. Those changes made the route into the goal more understandable and dependable without lowering the intellectual goal. The evidence was that adult language, pacing, modeling, and response options belong inside the explanation of participation. Gardner and Jorgensen deepen that lesson from different directions. First-hand accounts reveal harms that visible behavior counts can miss, while communicative competence depends on access to meaningful language throughout the day \citep{gardner2017firsthand,jorgensen2018competence}. Together, the readings require me to ask not only whether a strategy changes behavior, but whether it increases the learner's ability to communicate, set boundaries, participate in rigorous work, and be understood.

\reviewsection{The Media Informed My Judgment Without Replacing It}

\vspace{-0.2cm}
The Edutopia classroom-management session emphasizes structure, predictable routines, relationships, instructional clarity, and well-designed transitions \citep{edutopiaClassroom}. Its strongest contribution is that classroom management is partly lesson design: confusion, hidden sequences, and weak transitions can create conditions adults later describe as student behavior problems. The source also gave me practice applying EQUITAS to professional advice. Its interpretation of a multilingual learner moves quickly toward resistance and punishment without fully examining language access, culture, disability, trauma, school trust, or adult power. I can use its prevention-oriented ideas while refusing to inherit every interpretation. A source should inform my professional judgment. Respectful support requires the same discipline with research, media, checklists, and lived experience: identify what each contributes, name what it cannot establish, and keep missing voices visible.
